\documentclass[11pt]{book}

%-----------------------------------------%
%------------- PACKAGES ------------------%
%-----------------------------------------%
\usepackage[table]{xcolor}
\usepackage[colorlinks=true,linkcolor=blue,urlcolor=blue,citecolor=blue]{hyperref}
\usepackage[utf8x]{inputenc}
\usepackage{bm}	% bold in math mode \usepackage{bbm}
\usepackage{amsfonts,amsmath,amssymb,stmaryrd}
\usepackage{graphicx}

%-----------------------------------------%
%--------------- BIB ---------------------%
%-----------------------------------------%

\usepackage[numbers,square]{natbib}	% Natural bibliography with numbers and squares
\bibliographystyle{plain}		% Bib style

%-----------------------------------------%
%-------- New and Renew Commands ---------%
%-----------------------------------------%

%------------- Physics -------------------%
\renewcommand{\vec}[1]{\bm{#1}}
\newcommand{\bra}[1]{\langle#1|}
\newcommand{\ket}[1]{|#1\rangle}
\newcommand{\bkt}[2]{\left\langle#1|#2\right\rangle}
\newcommand{\comm}[2]{\left[#1,#2\right]}
\newcommand{\acomm}[2]{\left\{#1,#2\right\}}

\newcommand{\ad}{\hat{a}^\dagger}
\renewcommand{\a}{\hat{a}}
\renewcommand{\b}{\hat{b}}
\newcommand{\bd}{\hat{b}^\dagger}
\renewcommand{\d}{\hat{d}}
\newcommand{\dd}{\hat{d}^\dagger}
\renewcommand{\c}{\hat{c}}
\newcommand{\cd}{\hat{c}^\dagger}

\newcommand{\Svec}{\hat{\vec{S}}}
\newcommand{\Sz}{\hat{S}^z}
\newcommand{\Sp}{\hat{S}^+}
\newcommand{\Sm}{\hat{S}^-}
\renewcommand{\H}{\hat{\mathcal{H}}}
\renewcommand{\P}{\hat{\mathcal{P}}}

\newcommand{\ion}{\text{ion}}
\newcommand{\hc}{\text{h.c.}}
\newcommand{\Ry}{\text{Ry}}
\newcommand{\eff}{\text{eff}}


\newcommand{\pdv}[2]{\frac{\partial #1}{\partial #2}}
\newcommand{\tdv}[2]{\frac{\text{d} #1}{\text{d} #2}}

\newcommand{\intdV}{\int \text{d}^3 r}
%---------- non-Physics ------------------%
\newcommand{\blankpage}{
\newpage
\thispagestyle{empty}
\mbox{}
\newpage
}% Create a blankpage

%-----------------------------------------%
%--------------- GEOMETRY ----------------%
%-----------------------------------------%

\usepackage{geometry}
\geometry{
	hmargin=1.5cm, % Horizontal margin
	vmargin=1.75cm, % Vertical margin
	a4paper, % Paper size, change to letterpaper for US letter size
	%showframe, % Uncomment to show how the type block is set on the page -- typically for debugging
}
%-----------------------------------------%
%------------ MATH OPERATOR --------------%
%-----------------------------------------%


\DeclareMathOperator{\sinc}{sinc}
\DeclareMathOperator{\const}{const}
\DeclareMathOperator{\Tr}{Tr}   %trace
\DeclareMathOperator{\Dim}{dim}

%-----------------------------------------%
%------------- Title Page ----------------%
%-----------------------------------------%

\title{\Huge Condensed Matter Script}
\author{Petar \v{C}ubela}
\date{\today}

%-----------------------------------------%
%------------ Document -------------------%
%-----------------------------------------%
\begin{document}
% Create titlepage
\maketitle
\newpage
\tableofcontents
\newpage


%-----------------------------------------%
%-----------------------------------------%
%-----------------------------------------%
\chapter{Introduction}\label{chapter1}

Provide background information...

%-----------------------------------------%
\chapter{Separation of Lattice and Electron Dynamics}
%-----------------------------------------%
\section{The Theory of Everything}

Particle and String theorists search to find a better way to formulate the fundamental laws of physics. But suppose for a moment, we knew the fundamental ``theory of everything'' (TOE). Does this mean that the physics would stop existing as a basic science? Well, in condensed matter physics (CMP) we have a TOE, yet major discoveries continue to take place. The TOE in CMP is the many-particle Schroedinger equation
\begin{equation}
  i\hbar \pdv{}{t} \psi = \hat{H} \psi
\end{equation}
with the Hamiltonian 
\begin{equation}
  \label{TOE}
  \hat{H} = T_{\text{e}} + T_{\text{i}}+ V_{\text{e},\text{e}}+ V_{\text{i},\text{i}}+ V_{\text{i},\text{e}}
\end{equation}
where\\
\textcolor{red}{TODO: align the equations}
\begin{eqnarray}
  T_{\text{e}} = - \sum_{j=1}^{N_e} \frac{\hbar^2}{2 m_e} \nabla_j^2 \quad T_{\text{i}} = - \sum_{\ell=1}^{N_i} \frac{\hbar^2}{2 M_{\ell}} \nabla_{\ell}^2\\
  V_{\text{e},\text{e}} = \sum_{j,j^\prime =1}^{N_e} \frac{e^2}{\left|\vec{r}_j - \vec{r}_{j^\prime}\right|}\quad   V_{\text{i},\text{i}} = \sum_{\ell,\ell^\prime =1}^{N_i} Z_{\ell} Z_{\ell^\prime}\frac{e^2}{\left|\vec{R}_{\ell} - \vec{R}_{\ell^\prime}\right|} \\
  V_{\text{i},\text{e}} = -\sum_{j,\ell =1}^{N_e,N_i} \frac{Z_{\ell}e^2}{\left|\vec{r}_j - \vec{R}_{\ell}\right|}
\end{eqnarray}
Here $\vec{r}_j (\vec{R}_{\ell})$ refers to the coordinates of the $N_{\ell}(N_{\ion})$ electrons (ions) with mass $m_e(M_{\ell})$. $Z_{\ell}$ is the corresponding ion charge. The wave function depends on all coordinates \[\psi(\vec{r}_1,\dots, \vec{r}_{N_e}, \vec{R}_1, \dots,\vec{R}_{N_{\ion}})= \psi \left(\{\vec{r}_j\},\{\vec{R}_{\ell}\}\right) \] 
with the exeption of radiation effects and spin-orbit interactions (both can easily be included) all phenomena of CMP can be described by this Hamiltonian and the corresponding equation of motion.

We want to describe the coupled electrons and ions in a solid. It is clear that due to charge neutrality \[N_e = \sum_{\ell =1}^{N_{\ion}} Z_{\ell}.\]
Every atom in a solid consists of three parts: the positively charged heavy nucleus at the center, the light cloud of the many negatively charged core electrons which are tightly bound to the nucleus, and finally, the outer most few valence electrons. We refer to the nucleus with its core electrons as an ion.

\textcolor{red}{TODO: illustration of the ions and electrons}

Due to the positively charged background of ions which neutralize the "sea" of interacting electrons the self-interacting terms in the above Hamiltonian will be canceled: \[V_{ee}= \frac{1}{2} \sum_{j \neq i = 1 }^{N_e} \frac{e^2}{\left|\vec{r}_i - \vec{r}_j\right|} \quad V_{i,i}= \frac{1}{2} \sum_{\ell \neq k = 1 }^{N_{\ion}} \frac{Z_{\ell}Z_{k}e^2}{\left|\vec{R}_{\ell} - \vec{R}_k\right|}.\]

We will subsequently show that it is possible to trat the electron gas and the ionic motion separately. To be able to compare the relative magnitude of the terms in~\eqref{TOE}, introduce so-called \textit{atomic units}.

The characteristic length scale in a solid (or any condensed state of matter) is the interparticle spacing which is of the order of some \AA. The natural length scale in this regime is therefore the Bohr radius $a_B$, \[a_B = \frac{\hbar^2}{m_e e^2} = 0,053 nm\] and correspondingly we measure in units of \[E_0 = \frac{e^2}{a_B} = 2 \Ry = 27,2 eV.\]

Another useful length scale to consider is the interparticle spacing $ r_0 $  between electrons which is defined by the relation, \[ V = \frac{4}{3} \pi r_0^3 N_e.\] Using $ r_0 $ we can define a dimensionless quantity with which one can  distinguish the low and high-density regime in an interacting electron gas, namely: \[r_s := \frac{r_0}{a_B}.\] Using that the electron density $ n_e = N_e / V $ can be represented using the Fermi wave vector, $ n_e = k_F^3 / 3\pi^2 $ we can rewrite $ r_s $ as \[ r_s = \left( \frac{9\pi}{4} \right)^{1/3} \frac{1}{a_B k_F}.\]  

Using the atomic units we define a dimensionless vector $ \tilde{\vec{r}} := a_B^{-1} \vec{r} $ and rewrite the Hamiltonian~\eqref{TOE} as:
\begin{equation}
  \H = \frac{e^2}{2a_B} \left[- \sum_{j=1}^{N_e} \nabla^2_{\tilde{r}}- \sum_{\ell=1}^{N_{\ion}} \frac{m_e}{M_{\ell}}\nabla^2_{\tilde{R}} + \sum_{j\neq i =1 }^{N_e} \frac{1}{\left| \tilde{\vec{r}}_i - \tilde{\vec{r}}_j\right| } + \sum_{\ell \neq k =1 }^{N_{\ion}} \frac{Z_{\ell}Z_k}{\left| \tilde{\vec{R}}_{\ell} - \tilde{\vec{R}}_k\right| } -2\sum_{j\neq \ell =1 }^{N_e, N_{\ion}} \frac{Z_{\ell}}{\left| \tilde{\vec{r}}_j - \tilde{\vec{R}}_{\ell}\right| } \right]
\end{equation}
Apparently, the Hamiltonian depends in the new units only on the ionic charge numbers $Z_{\ell}$ and on the mass ratio $ \frac{m_e}{M_\ell} $. In particular, the relative contribution of the kinetic energy of the ions compared to the other terms in the Hamiltonian is exactly this factor $ \frac{m_e}{M_{\ell}} $. Hence, the kinetic energy of the ions is smaller about this ratio. Due to the magnitude of $ \frac{m_e}{M_{\ell}} $ of around $10^{-4} - 10^{-5}$ it is appropriate to expand the theory with respect to this parameter and hence the kinetic energy of the ions.

Hence, in the lowest order of perturbation theory it will be possible to assume that the ions do not move and the electrons therefore propagate in the rigid lattice of charge neutralizing ions.

The smallness of the ration $ \frac{m_e}{M_{\ell}} $ is the reason to assume that the motion of the electrons and lattice vibrations decouple. Corrections are treated in perturbation theory. 

In the next section we will show how the electronic and lattice degrees of freedom decouple using an perturbative approximation.


%-----------------------------------------%
\section{Born-Oppenheimer Approximation}

Following the arguments of the previous chapter we will treat the ionic kinetic energy as a small perturbation to the other terms.
\begin{equation}
  \H = T_{\ell} + V_{e,e} + V_{e,\ion} + V_{\ion, \ion} + T_{\ion} \equiv \H_0 + T_{\ion}
\end{equation}
We want to analyze the Schroedinger equation
\begin{equation}
  \H \Psi(\vec{r},\vec{R}) = E \Psi(\vec{r},\vec{R})
\end{equation}
Assume that eh eigenvalue problem of the unperturbed system is solved and the solution known:
\begin{equation}
  \H_0 \psi_{\alpha}\left( \{\vec{r}_i\}, \{ \vec{r}_j\} \right) = \varepsilon_{\alpha}\left( \{\vec{R}_j\}\right) \psi_{\alpha}\left( \{\vec{r}_i\}, \{ \vec{r}_j\} \right)
\end{equation}
The ion positions $ \vec{R}_j $  enter this equation only as parameters because the positions of the ions are not dynamical in this problem. Hence, the electronic energy spectrum depends only on the rigid positions of the charge neutralizing lattice. Note that the ion-ion interaction term is only an additive constant for the electronic spectrum, so we can ignore it in $ \H_0 $ and include it later when including the perturbation. The eigenfunctions $ \psi_{\alpha}\left( \{\vec{r}_i\}, \{ \vec{r}_j\} \right) $ build a complete set of orthonormal basis functions for a given set of ion positions. 

Therefore, the general eigenfunction of the eigenvalue problem 
\begin{equation}
  \H \Psi \left( \vec{r}, \vec{R} \right) = E \Psi \left( \vec{r}, \vec{R} \right)
\end{equation}
can be expanded using the electronic wave functions $ \psi_{\alpha} \left( \vec{r}_i, \vec{R}_j \right) $ for a given set of ion positions:
\begin{equation}
  \Psi \left( \vec{r}, \vec{R} \right) = \sum_{\alpha} \chi_{\alpha}\left( \vec{R} \right) \psi_{\alpha} \left( \vec{r}, \vec{R} \right)
\end{equation}
Inserting this expansion into the eigenvalue equation for $ \H $ leads to 

\begin{eqnarray}
  \H \Psi \left( \vec{r}, \vec{R} \right) = \left(\H_0 + V_{\ion, \ion} + T_{\ion} \right) \sum_{\alpha} \chi_{\alpha}\left( \vec{R} \right) \psi_{\alpha} \left( \vec{r}, \vec{R} \right) \\ 
  = \sum_{\alpha} \left( \varepsilon_{\alpha}\left( \vec{R} \right) + V_{\ion, \ion} + T_{\ion} \right)\chi_{\alpha}\left( \vec{R} \right) \psi_{\alpha} \left( \vec{r}, \vec{R} \right)
\end{eqnarray}
which can be rewritten into 
\begin{equation}
  \sum_{\alpha} \left( \varepsilon_{\alpha}\left( \vec{R} \right) + V_{\ion,\ion} + T_{\ion} - E \right)\chi_{\alpha}\left( \vec{R} \right) \psi_{\alpha} \left( \vec{r}, \vec{R} \right) = 0.
\end{equation}
We multiply the this equation by $ \psi_{\beta}^* \left( \vec{r}, \vec{R} \right) $  and integrate over $ \vec{r} $ (using $ \intdV  \psi_{\beta}^* \left( \vec{r}, \vec{R} \right) \psi_{\alpha} \left( \vec{r}, \vec{R} \right) = \delta_{\alpha, \beta}$ ) which leads us to 
\begin{equation}
  \left( \varepsilon_{\beta}\left( \vec{R}\right) + V_{\ion, \ion} + T_{\ion} -E \right) \chi_{\beta} \left( \vec{R} \right) + \sum_{\alpha} A_{ \beta, \alpha}\left( \vec{R} \right) \chi_{\alpha}\left( \vec{R} \right) = 0
\end{equation}
where 
\begin{equation}
  A_{ \beta, \alpha}\left( \vec{R} \right) = \intdV \sum_{\ell=1}^{N_{\ion}} \frac{1}{2M_{\ell}} \psi_{\beta}^* \left( \vec{r}, \vec{R} \right) \left( \nabla^2_{\vec{R}_{\ell}}\psi_{\alpha} \left( \vec{r}, \vec{R} \right) + 2 \nabla_{\vec{R}_{\ell}}\psi_{\alpha} \left( \vec{r}, \vec{R} \right)\nabla_{\vec{R}_{\ell}} \right)
\end{equation}
$ A_{ \beta, \alpha}\left( \vec{R} \right) $ is an overlap matrix element between different electronic states; if it is possible to ignore this term we arrive at 
\begin{equation}
  \left( T_{\ion} + \varepsilon_{\beta}\left( \vec{R} \right) +V_{\ion,\ion}\right) \chi_{\beta} \left( \vec{R} \right) = E \chi_{\beta} \left( \vec{R} \right)
\end{equation}
which is an eigenvalue equation for the ionic degrees of freedom only. The effective potential here is 
\begin{equation}
  V_{\eff} \left( \vec{R} \right) = \varepsilon_{\beta} \left( \vec{R} \right) + V_{\ion,\ion}
\end{equation}
and thus the electronic energy spectrum determines the spectrum of the lattice dynamics. All effects of chemical bonds of the ions and the energetically favored lattice structures are encoded in the minimized electronic energy $ \varepsilon_{\beta}\left( \vec{R} \right) $ for some special postion $ \vec{R}_0 $. 






\clearpage
\newpage
%-----------------------------------------%
\chapter{Crystal Structure}\label{chapter2}



%-----------------------------------------%
\section{Bravais Lattice}\label{section2}
There are two equivalent definitions for a Bravais lattice:
\begin{itemize}
  \item A Bravais lattice is an infinite array of discrete points with an arrangement and an orientation that appears exactly the same from whichever of the points the array is viewed.
  \item A $d$-dimensional Bravais lattice consists of all points with position vector 
    \begin{equation}
      \vec{R} = \sum_{i=1}^{d} n_i \vec{a}_i\quad n_i \in \mathbb{Z}
    \end{equation}
      where $\{\vec{a}_i\}$ are linearly independent vectors. They do not have to be orthogonal. They are called \textit{lattice basis vectors} and are not unique. In general there an infinite number of lattice basis vectors leading to the same Bravais lattice.
\end{itemize}


\clearpage
\newpage

%-----------------------------------------%
\chapter{Lattice Vibrations (Phonons)}\label{chapter3}

\clearpage
\newpage

%-----------------------------------------%
\section{Collective Oscillations of Crystal Structures}


\clearpage
\newpage

%-----------------------------------------%
\chapter{Numerics for lattice systems and phonons}

%-----------------------------------------%
\section{Introduction}

Goal:\\ 
Python script which simulates oscillations of the lattice (with pygame). This should work for arbitrary crystal structures (in 2D and 3D). Calculate and plot the dispersion of phonon modes.

Milestones:
\begin{itemize}
  \item Library to solve second order initial value problems numerically
  \item class to plot lattice structures
  \item library to plot dispersion relations along symmetry points
  \item (symbolic python to get ODE from potential)
\end{itemize}

%-----------------------------------------%
\chapter{Non-Interacting Electrons in the Solid}

%-----------------------------------------%
\section{Bloch's Theorem}\label{section4.1}
In a real solid the electrons are subjected to a periodic potential $V(\vec{r}) = V(\vec{r} + \vec{R})$. Here $\vec{R}$ are the lattice vector of the underlying crystal structure. Thus, the Hamiltonian operator 
\begin{equation}
  \H = \sum_{i} \left(\frac{\vec{p}^2_i}{2m} + V(\vec{r}_i) \right)
\end{equation}
has a discrete translational symmetry. 
\clearpage
\newpage


%-----------------------------------------%
\chapter{Electron-Electron Interactions}\label{chapter5}

%-----------------------------------------%
\chapter{Electron-Phonon Interactions}\label{chapter4}

%-----------------------------------------%
\chapter{Doped Spin Chain in staggered Field}\label{chapter5}

%-----------------------------------------%
\section{The t-J Model}\label{section2.1}
The $t-J$ model is given by:
\begin{equation}
  \H = J \sum_j \Svec_{j+1}\Svec_j - h\sum_j {(-1)}^{j} \Sz_j
  -t \sum_{j\sigma} \P \left(\cd_{j+1\sigma}\c_j + \hc\right)\P
\end{equation}
This Hamiltonian posses a continuous $SU(2)$ symmetry.

In general there is no spontaneous breaking of a continuous symmetry in 1D.
\textcolor{red}{Add more context to the above statement!. Start with the t-J model and explain that one includes the staggered field such that the symmetry is explicitly broken.} 

Thus, the model of interest is the $t-J$ model in a staggered field,
\begin{equation}
  \H = J \sum_j \Svec_{j+1}\Svec_j - h\sum_j {(-1)}^{j} \Sz_j
  -t \sum_{j\sigma} \P \left(\cd_{j+1\sigma}\c_j + \hc\right)\P
\end{equation}
where we assume that one hole is doped into the system. I.e. \[ \sum_{j\sigma} \cd_{j\sigma}\c_{j\sigma} = N_e -1\] where $N_e$ is the number of electrons into the system. Due to the fact that we consider the system slightly away from half filling, the electron number $N_e$ equals the number of lattice sites $N$.


%-----------------------------------------%
\section{Ising Model in a staggered Field}\label{section2.2}
We first look at the physics of the simple Ising model in staggered magnetic field:
\begin{equation}
  \H = J_z \sum_j \Sz_{j+1}\Sz_{j} - h \sum_j {(-1)}^{j}\Sz_j.
\end{equation}
The Ising model shows ordering at $T=0$ and establishes one of the two states \[ \ket{\dots\uparrow\downarrow\uparrow\downarrow\uparrow\dots}\quad\ket{\dots\downarrow\uparrow\downarrow\uparrow\downarrow\dots}.\] The system chooses one of the two degenerate states depending on the sign of $h$. In the following we will set without loss of generality $h>0$! Thus, for a non-periodic chain the ground state of the system is 
\begin{equation}
  \label{neel_state}
  \ket{\text{N\'eel}} := \ket{\downarrow\uparrow\downarrow\dots}, \quad E_{\text{GS}} =- J_z \frac{1}{4} \left(N-1\right) -N \frac{h}{2}.
\end{equation}
The excitations in this model are dispersionless but gapped Magnon excitations with the excitation energy $\varepsilon = J_z + h$.

For this system one is able to compute the canonical partition function $Z$ exactly. 
\textcolor{red}{The canonical partition is given by 
\begin{equation}
  Z = \Tr\left(e^{-\beta \hat{\mathcal{H}} }\right)
\end{equation}
where $\beta= 1/k_B T$ is the inverse temperature and $k_B$ corresponds to the Boltzmann constant.
}

%-----------------------------------------%
\section{From Ising to the full Model}\label{section2.3}
We ask the question if it is more practical to first extend the Ising model to the $t-J_z-h$ model or first to the XXZ-model in a staggered field. 
\begin{itemize}
  \item XXZ-$h$ $=$ AFM ground state $+$ spin fluctuations represented by magnon excitations 
  \item $t-J_z-h$ $=$ AFM ground state $+$ charge fluctuations represented by a chargon bound to a localized spinon
\end{itemize}

\textcolor{red}{Include / Draw the figure in my note book. Page 6.}

We comment on the chosen parameters in the model. They are chosen such that 
\begin{equation}
  t \ll J_z, J_\perp, \quad J_z \le J_\perp, h \text{ arbitrary}
\end{equation}
Due to the large value of $t$ the theory is a strong coupling theory and thus the charge and spin degrees of freedom can be separated by a Born-Oppenheimer approximation.


%-----------------------------------------%
\section{XXZ Model in a staggered field\label{section2.4}}

Here we next consider the XXZ-model in a staggered field,
\begin{equation}
  \H_{\text{XXZ}} = J_z \sum_j \Sz_{j+1}\Sz_{j} + \frac{J_\perp}{2} \sum_{j} \left( \Sp_{j+1}\Sm_j+ \hc\right) - h \sum_j {(-1)}^{j} \Sz_j.
\end{equation}
The ground state of this model is an AFM state; the ``flip-flop'' terms lead to fluctuations of the spins above the perfect N\'eel state.
The spin fluctuations lead to coherent spin waves which can be represented by magnon excitations. The best way to introduce them is the so-called Holstein-Primakoff transformation. Due to this transformation deviations from the perfect N\'eel state~\eqref{neel_state} are represented by quantum excitations. Importantly, the transformation has to be done carefully because the transformation differs depending on the spin state on a lattice site. It turns out that for the above Hamiltonian one can expand the system around the perfect N\'eel state, $ \ket{\dots\uparrow\downarrow\uparrow\downarrow\uparrow\dots} $. 
The Holstein-Primakoff transformation can in approximated form be written as 
\begin{equation}
  \Sz_j = \left( -1\right)^{-1}\left( S - \a^\dagger_j \a_j \right)
\end{equation}

%-----------------------------------------%
\section{Doped hole in the Ising model in a staggered Field}\label{section2.5}


\clearpage
\newpage
%%%%%%%%%%%%%%%%%%%%%%%%%%%%%%%%%%%%%%%%%%%%%%%%%%%%%%%%%%%%%%%%%%%%%%%%%%%%%%%%%
%-----------------------------------------%
\chapter{MixD doped AFM}\label{chapter6}



\clearpage
\newpage
%%%%%%%%%%%%%%%%%%%%%%%%%%%%%%%%%%%%%%%%%%%%%%%%%%%%%%%%%%%%%%%%%%%%%%%%%%%%%%%%%
\chapter{2D doped AFM}\label{chapter7}


\clearpage
\newpage


%-----------------------------------------%
%-------------- Appendix -----------------%
%-----------------------------------------%
\appendix

\setlength\parindent{0pt}

%-----------------------------------------%
\chapter{Useful Formulas}

%-----------------------------------------%
\section{Baker-Campbell-Hausdorff}

\begin{equation}
  e^{X} Y e^{-X} = \sum_k \dfrac{1}{k!} \comm{X}{Y}_k
\end{equation}
with \[\comm{X}{Y}_k = \comm{X}{\comm{X}{Y}_{k-1}},\quad\quad
\comm{X}{Y}_0.\]


%-----------------------------------------%
\section{Trigonometrics}

%-----------------------------------------%
\subsection{Double angle}

\begin{itemize}
  \item $ 2\sin^2(x) = 1 - \cos(2x) $ 
\end{itemize}


%-----------------------------------------%
\section{Gamma Function}

The so-called \textit{Gamma function} is defined as 
\begin{equation}
  \Gamma(z) = \int_0^{\infty} \text{dt}\: t^{z-1} e^{-t}, \quad \Re(z) > 0.
\end{equation}

It has the following properties:
\begin{itemize}
  \item $ \Gamma(n+1) = n!, \quad n= 0, 1, 2, \dots $ 
  \item $ \Gamma(1) = 0! = 1 $ 
  \item $ \Gamma(1/2) = \sqrt{\pi} $ 
  \item $ \Gamma(z+1) = z \Gamma(z) $ 
  \item $ \Gamma(z)\Gamma(z+1) = \frac{\pi}{\sin(\pi z)} $ 
\end{itemize}
The Gamma function is generalization of the faculty operation for $ z \in \mathbb{C} $. It is widely used in different areas of physics.

One important case is the solution of the integral \[I_{\nu} = \int_0^{\infty} \text{dy} \frac{y^{\nu-1}}{e^y -1}\] which appears in the treatment of Bose gases. Its solution is given by: \[ I_{\nu} = \Gamma(\nu) \zeta(\nu),\] where $ \zeta(\nu) = \sum_{n=1}^{\infty} \frac{1}{n^{\nu}} $ is the famous Riemann Zeta function.

\clearpage
\newpage

%-----------------------------------------%
\chapter{Pauli Matrices}
The Pauli matrices are defined as:
\begin{equation}
  \sigma_0=\begin{pmatrix}
                1 & 0\\
                0 & 1
            \end{pmatrix}
  \quad
  \sigma_x= \begin{pmatrix}
                0 & 1\\
                1 & 0
            \end{pmatrix}
  \quad
  \sigma_y= \begin{pmatrix}
                0 & -i\\
                i & 0
            \end{pmatrix}
  \quad
  \sigma_z = \begin{pmatrix}
                1 & 0\\
                0 & -1
            \end{pmatrix}
\end{equation}
Strictly speaking the matrix $\sigma_0$ is not a Pauli matrices but it will be useful to give it this name.

For convenience we define the symbolic vector
\begin{equation}
      \vec{\sigma} = \left(\sigma_x, \sigma_y,\sigma_z \right)^T.
\end{equation}


We record some properties of the Pauli matrices. First,
\begin{eqnarray}
  \sigma_i^2 =1  \\
  \sigma_i \sigma_j + \sigma_j \sigma_i = 0, \quad \text{for } i\neq j.
\end{eqnarray}
where the right hand side of the first identity is to; be understood as the $2\times 2$ identity matrix. These two relations are equivalent to the anti-commutation relations 
\begin{equation}
  \acomm{\sigma_i}{\sigma_j} = 2 \delta_{ij}.
\end{equation}
We also have the commutation relation 
\begin{equation}
  \comm{\sigma_i}{\sigma_j} = 2i \varepsilon_{ijk} \sigma_k,
\end{equation}
which is the $2\times 2$ matrix realization of the angular momentum commutation relation. Notice also that 
\begin{eqnarray}
  \sigma_i^\dagger = \sigma_i,\\
  \det(\sigma_i) = -1\\
  \Tr(\sigma_i) = 0,
\end{eqnarray}
which means that the Pauli matrices are hermitian and traceless. In addition the Pauli matrices are linear independent. Therefore they build a span for the space of hermitian matrices. 

Every hermitian $2\times 2$ matrices $A$ can be written as the linear combination,
 \begin{equation}
   A = a_0 \sigma_0 + \vec{a}\cdot \vec{\sigma}
 \end{equation}
where we included the identity matrix $\sigma_0$.
Additionally, the Pauli matrices are the generators for special unitary matrices, $SU(2)$.
The following exponential describes rotations in the $2$-component spinor space:
        \begin{equation}
              \exp(i \vartheta \: \vec{n} \cdot \vec{\sigma} ) = \cos(\vartheta) \sigma_0 + i \vec{n} \cdot \vec{\sigma} \sin(\vartheta)
        \end{equation}
      and is a very useful relation in computations.

\clearpage
\newpage

%%%%%%%%%%%%%%%%%%%%%%%%%%%%%%%%%%%%%%%%%%%%%%%%%%%%%%%%%%%%%%%%%%%%%%%%%%%%%%%%%%%%%%%%%%%%%%
%%%%%%%%%%%%%%%%%%%%%%%%%%%%%%%%%%%%%%%%%%%%%%%%%%%%%%%%%%%%%%%%%%%%%%%%%%%%%%%%%%%%%%%%%%%%%%
%-----------------------------------------%
\chapter{Multi Coherent States}



\clearpage
\newpage

%-----------------------------------------%
\chapter{Bogoliubov transformation}

\section{Single Mode}

In several quite interesting physical problems we encounter a Hamiltonian of the form
\begin{equation}
  \H = \sum_q A_q \ad_q \a_q + \frac{1}{2} \sum_q \left(B_q \a_q \a_{-q} + \hc\right)
\end{equation}
where  $ A_q = A^*_q $ and $ B_q = B_{-q} $ follows by symmetry considerations. The quantum numbers $ q $ are momenta in $ d $-dimensions.   These coefficients differ depending on the model in question. This Hamiltonian is important for excitation in magnetic systems, the excitation spectrum in a superfluid and for the excitations in the BCS theory of superconductivity.

This Hamiltonian can be diagonalized by introducing new bosonic annihilation and creation operators,
\begin{eqnarray}
  \b_{q} = u_{q} \a_q + v_q \ad_{-q}\\
  \bd_{q} = u^*_{q} \ad_q + v^*_q \a_{-q}.
\end{eqnarray}
Consequently, the transformation has to be canonical which is the case for 
\begin{equation}
  \left| u_q \right|^2 - \left| v_q \right|^2 \overset{!}{=} 1
\end{equation}
This form suggests to use the relation \[ \cosh^2(x) - \sinh^2(x) = 1. \] Therefore we can choose

\begin{eqnarray}
  u_q = e^{i\varphi_1} \cosh(\vartheta_q)\\
  v_q = e^{i\varphi_2} \sinh(\vartheta_q)
\end{eqnarray}

For simplicity and not-knowledge we set $ \varphi_1 = \varphi_2 = 0 $. 

By introducing the new bosonic modes into the Hamiltonian and demanding that it is diagonal we get the following relations for the so-called Bogoliubov coefficients $ u $ and $ v $,
\begin{eqnarray}
  u_q^2 + v_q^2 = \cosh(2\vartheta_q) = \frac{A_q}{\omega_q}, \quad 2u_q v_q = \sinh(2\vartheta_q) = \frac{B_q}{\omega_q}
\end{eqnarray}
where 
\begin{equation}
   \omega_q = \sqrt{A_q^2 - B_q^2}
\end{equation}
is the spin-wave dispersion of the new bosonic modes:
\begin{equation}
  \H = \sum_q \omega_q \left(\bd_q\b_q + \frac{1}{2}\right) - \frac{1}{2} \sum_{q} A_q.
\end{equation}


\section{Multi Mode}
Here we consider a generic bosonic quadratic Hamiltonian operator,
\begin{equation}
  \H = \sum_{pq} A_{pq} \bd_p \b_q + \dfrac{1}{2} \sum_{pq} \left( B_{pq} \bd_p \bd_q +\hc \right)
\end{equation}
where the annihilation and creation operators fulfill the standard bosonic commutation relations,
\begin{equation}
  \comm{\b_p}{\bd_q} = \delta_{p,q} \quad \text{and} \quad \comm{\b^{(\dagger)}_p}{\b^{(\dagger)}_q} = 0.
\end{equation}

We want to study the diagonalizability of $ \H $ . The matrices in the Hamiltonian fulfill \[A = A^\dagger \quad{\text{and}}\quad  B = B^T\] by symmetry considerations.\\
Because of the fact that the Hamiltonian is quadratic, we expect it to take a form
\begin{equation}
    \H = E_0 + \sum_l w_l \dd_l\d_l
\end{equation}
in its diagonal basis.\\
This is achieved by applying a generalized Bogoliubov transformation
\begin{eqnarray}
    \d_l = \sum_q \left[U_{l,q}\b_q -V_{l,q}\bd_q\right]\\
    \dd_l = \sum_q \left[U_{l,q}\bd_q - V_{l,q}\b_q \right]
\end{eqnarray}
and search $U,V$ such that
\begin{equation}
  \comm{\H}{\dd_l} = w_l \dd_l
\end{equation}
is fulfilled. The insertion of the expressions for the new bosonic operators $ \d_{l} $ into the commutator relation yields the following eigenvalue equation:
\begin{equation}
  \begin{pmatrix}
    A & B \\ 
    -B^* & -A*
  \end{pmatrix}
  \begin{pmatrix}
    U \\ V^*
  \end{pmatrix}_{\ell} 
  = w_{\ell} 
  \begin{pmatrix}
    U \\ V^*
  \end{pmatrix}_{\ell}
\end{equation}
Interestingly, the matrix on the left hand site of this equation is not hermitian. Therefore, there is no unitary transformation to diagonalize it.



\clearpage
\newpage
%-----------------------------------------%
\chapter{Second Qunatization}





\clearpage
\newpage

\bibliography{/home/petar/work/tex/physics.bib}

\end{document}
